\documentclass[10pt,a4paper]{article}
\usepackage[margin=1.7cm]{geometry}
\usepackage{enumitem}\usepackage{titlesec}\usepackage[hidelinks]{hyperref}\usepackage{parskip}
\titleformat{\section}{\normalfont\large\bfseries}{\thesection}{0.6em}{}
\titlespacing*{\section}{0pt}{6pt}{2pt}
\setlist[itemize]{topsep=2pt,itemsep=1pt,parsep=0pt,leftmargin=1.3em}
\pagestyle{empty}
\begin{document}
\begin{center}
{\Large\bfseries Motor Speed Control with a PI Controller and Anti-Windup}\\[2pt]
{\large Lab 3 \textendash{} digital control of the T-Bot motors (and Lab 4 line following)}\\[4pt]
{\small Embedded Real-Time Control laboratory \textendash{} M.Sc. Control Systems Engineering, University of Padova}\\
{\small Author: Kiavash Rabbani Soudabi (ID 2141822) \textendash{} Course held by Prof.\ Alberto Morato \textendash{} 2025}
\end{center}

\section*{Purpose}
Implement a PID controller for the T-Bot drive motors: measure the wheel speed from the quadrature encoders, close
the loop around the motor voltage, and command a reference speed from the robot keypad, with telemetry for tuning in
MATLAB.

\section*{What I implemented}
\begin{itemize}
  \item \textbf{Speed measurement:} I read the TIM3/TIM4 encoder counters and computed the per-period count
        difference, explicitly handling counter overflow and underflow so that the direction and the magnitude stay
        correct; counts are converted to rad/s with the encoder/gear scaling.
  \item \textbf{Digital PI controller} executed in the timer period-elapsed callback at a fixed sampling time of
        \texttt{TS = 10\,ms}: the error is evaluated in rad/s, with the proportional and integral terms
        $P=0.871$, $I=7.5563$, and the integral action updated as $I \mathrel{+}= I\,T_s\,e$.
  \item \textbf{Anti-windup:} conditional integration -- if the output saturates \emph{and} the error would worsen the
        saturation, the integral step is removed, which keeps the integrator usable after long saturation episodes.
  \item \textbf{Output conditioning and drive:} the control signal (volts) is limited to $\pm6$\,V, scaled to PWM duty
        against the 12\,V battery and the timer ARR value, and applied with two drive modes --
        \texttt{ForwardCoast} (single-channel PWM) and \texttt{ForwardBreak} (complementary PWM for braking).
  \item \textbf{Operator interface:} the reference speed is entered on the keypad (up to three digits followed by
        \texttt{\#}, decoded with \texttt{atoi}), and the keypad decoding/handling is interrupt- and buffer-based.
  \item \textbf{Telemetry:} the control signal and the measured quantities are streamed through the course datalogger
        for plotting and tuning in MATLAB.
  \item \textbf{Lab 4 (line following)} reuses this exact motor-control base and adds the line sensor handling --
        storing the sensor pattern (\texttt{StoreBinary}) and computing the lateral error used by the steering law.
        Lab 4 has no separate report; its implementation lives in the same code base.
\end{itemize}

\section*{Results}
The speed loop tracks the commanded reference on both motors, with the measured step response and the control signal
recorded by the datalogger and analysed in MATLAB. The report documents the controller structure, the gains, the
anti-windup logic and the measured tracking behaviour. The same base was extended for the autonomous line-following
task of Lab 4.

\section*{Repository contents}
\texttt{report/Embedded-Report-lab3.pdf} (submitted report) $\cdot$ \texttt{src/user\_code\_main.c} (USER CODE
sections of my \texttt{main.c} -- speed measurement, PI controller with anti-windup, drive functions, keypad
command handling; the course framework around them is not redistributed) $\cdot$ \texttt{docs/} (this overview and
report pages).

\section*{Skills demonstrated}
Digital control implementation (PI, fixed sampling time, anti-windup), quadrature encoder interfacing and
overflow-safe speed estimation, PWM motor drives including braking mode, embedded real-time scheduling (timer
interrupt control loop), measurement, logging and tuning with MATLAB, bare-metal C on ARM Cortex-M.

\end{document}
